% Opening paragraph of Chapter 7, last sentence. Reason: the Conclusions ("Both halves of this thesis lead to the same advice") give the same advice.
Before choosing a method, we therefore ask whether the labels were pooled from many laboratories, measured in one controlled experiment or taken from an unknown source.

% Summary, paragraph on Chapter 6. Reason: the same result is asked for in "Simple Models and Small Feature Sets", paragraph 2 (the "fewer features" paragraph), which keeps it.
Nearest-neighbor prediction: \textbf{[RESULT NEEDED]}.

% Simple Models and Small Feature Sets, paragraph 1. Reason: state the Chapter 3 results in words and point to Chapter 3 (removed words marked; the rest of the sentence stays).
NGBoost ranked sixth to ninth on hERG $K_i$, while

% Simple Models and Small Feature Sets, paragraph 4. Reason: the Song2022 claim is stated in Section sec:lnl_motivation (Chapter 4) and Section sec:lnl_background (Chapter 2); the Grinsztajn2022 claim is stated in "Models and Metrics", Chapter 5. Replaced by short clauses with back-references.
Most were built and tested on large image datasets with flipped class labels \citep{Song2022}, whereas
On medium-sized tabular data, tree models still outperform deep networks, partly because uninformative features hurt networks more \citep{Grinsztajn2022}.

